%==========================
% Chapter 1 — Introduction
%==========================

\section{Introduction}

The rapid proliferation of Electric Vehicles (EVs) in Indian cities has created an
urgent need for scientifically planned public charging infrastructure. Inadequately
placed stations lead to range anxiety, long queuing times, and underutilised
investments --- outcomes that collectively erode user confidence and slow the
broader EV adoption transition. The city of Indore, a major commercial hub of
Madhya Pradesh, presents a challenging yet representative context for this problem,
given its diverse ward-level demography, heterogeneous land-use patterns, and
an active municipal push toward sustainable urban mobility under the Smart Cities
Mission. As of early 2024, the city had fewer than 50 publicly accessible EV
charging points --- a significant deficit relative to its projected EV fleet size
under any realistic adoption scenario.

Designing an optimized EVCS network is fundamentally a combinatorial
multi-objective problem. A planner must simultaneously minimize capital
expenditure, maximize spatial demand coverage, maximize financial return on
investment, and minimize customer queue waiting times --- objectives that are
inherently conflicting. No single solution dominates all others across every
objective; instead, a Pareto-optimal set of trade-off solutions exists, which
is the domain of multi-objective evolutionary algorithms (MOEAs). This thesis
addresses this problem through a comprehensive four-objective optimization
framework applied across the 85~administrative wards of Indore, comparing
two state-of-the-art evolutionary algorithms --- Non-dominated Sorting Genetic
Algorithm~II (NSGA-II) and Multi-Objective Particle Swarm Optimization
(MOPSO) --- and validating results through analytical queueing models and
Monte~Carlo stochastic demand simulation.

\subsection{Overview of the Proposed Framework}
\label{sec:intro_overview}

The global transition toward sustainable transportation has led to a rapid
increase in EV adoption. In India, and particularly in fast-growing cities
such as Indore, this trend is reinforced by government incentives under the
FAME~II scheme, environmental regulations, and rising public awareness.
However, large-scale EV uptake is fundamentally constrained by the
availability, accessibility, and reliability of the charging infrastructure.
Poorly planned EVCS networks can lead to range anxiety, long waiting times,
uneven spatial coverage, and overloading of the power grid.

Designing an optimized EVCS network is a complex, multi-objective task.
It involves simultaneously deciding:
\begin{itemize}
  \item which candidate locations should host charging stations;
  \item how many chargers and what capacity should be installed at each
        selected site;
  \item how charging prices should be set across the network under
        Time-of-Use (ToU) demand patterns;
  \item how to ensure that grid constraints, blocking probability thresholds,
        and user service quality requirements are simultaneously satisfied.
\end{itemize}
These decisions must be made under spatially heterogeneous demand,
limited grid capacity, stochastic demand uncertainty, and evolving
population and EV growth over a multi-year horizon.

To address this challenge, the thesis proposes a unified framework with
five tightly coupled components:
\begin{enumerate}
  \item \textbf{Four-objective problem formulation:} The EVCS placement
        problem is formulated as a binary combinatorial optimisation problem
        minimizing total capital expenditure~($f_1$) and mean queue waiting
        time~($f_4$), while maximizing spatial demand coverage~($f_2$) and
        net return on investment~($f_3$), over a pool of $30+$ heterogeneous
        candidate venues across the 85~wards of Indore.

  \item \textbf{NSGA-II and MOPSO comparative optimization:} Both
        algorithms are independently applied to generate Pareto-optimal
        station deployment plans. NSGA-II employs fast non-dominated sorting
        with crowding distance selection, while MOPSO uses an external archive
        with crowding-based leader selection. Algorithm quality is formally
        compared using Hypervolume~(HV), Generational Distance~(GD), and
        Spread performance indicators.

  \item \textbf{Monte~Carlo stochastic demand modelling:} Demand is
        projected using a logistic EV adoption curve anchored to CAGR-based
        ward-level demographic forecasts, and further overlaid with
        100-scenario Monte~Carlo simulation (noise fraction $v = 0.15$) to
        ensure that selected station configurations are robust to stochastic
        demand fluctuations.

  \item \textbf{M/M/c/K queueing model with Time-of-Use pricing:}
        Service quality at each station is captured by an M/M/c/K
        finite-capacity queueing model with $K = c + 8$ physical waiting
        bays, evaluated across three ToU demand periods --- Peak, Normal,
        and Idle. A blocking probability constraint ($P_{\text{block}} > 10\%$)
        is embedded directly into the evolutionary evaluation loop as a
        hard penalty, ensuring that no selected station configuration
        becomes queue-infeasible under peak demand.

  \item \textbf{Sensitivity analysis and scalability benchmarking:}
        The robustness of the optimal solutions is examined with respect
        to EV growth rate, electricity tariff, installation cost factor,
        and charger service rate. Scalability benchmarks confirm polynomial
        runtime growth for both algorithms, validating applicability to
        larger metropolitan deployments.
\end{enumerate}

Applied to Indore, the framework produces well-converged Pareto
fronts spanning the full coverage range. Across the union of the
NSGA-II and MOPSO archives, capital expenditures range from
\rupee39\,L for the smallest viable layout to \rupee411\,L for full
(100\,\%) ward coverage, with a marginal cost of approximately
\rupee6--7\,L per percentage point of coverage. Mean queue waiting
times across all Pareto-optimal layouts lie in the range
0.13--5.6~minutes, comfortably within the 20-minute service-level
target. Final hypervolume values are 0.9967 for NSGA-II and 1.0434
for MOPSO, indicating that the two metaheuristics produce
complementary fronts. The end-to-end pipeline (forecasting,
queueing, optimisation, evaluation, and visualisation) completes in
under four minutes of wall-clock time on a commodity laptop.

\section{Background}
\label{sec:intro_background}

The expansion of EV infrastructure in Indian cities requires optimization
techniques capable of handling multiple layers of complexity. Unlike
conventional fuel stations, EVCSs interact directly with the power
distribution network, making their planning and operation depend
simultaneously on transportation patterns, demographic trends, grid
constraints, and financial viability.

\subsection{EV Adoption and Urban Charging Needs}

India ranks as the world's third largest automobile market and has
set an ambitious target of 30\% EV penetration by 2030 under the
National Electric Mobility Mission Plan (NEMMP). As of 2024, India
had registered over four million EVs, with two-wheelers
constituting the dominant segment. For Indore specifically, the
2011 census ward populations are projected forward at a compound
annual growth rate of $r = 1.8\%$, consistent with the 2001--2011
municipal trend, and the EV penetration fraction $\alpha(t)$ is
modelled as a logistic curve with inflection at $t_0 = 2030$
(aligned with the NEMMP target year) and steepness $a = 0.3$.
Under this projection the city-wide EV fleet is expected to grow by
roughly an order of magnitude between 2025 and 2035, putting clear
upward pressure on the public charging backbone.

Such growth patterns imply that EVCS networks cannot be designed solely
for current demand. They must be robust to substantial increases in both
the number of EVs and their spatial distribution, motivating the use of
logistic adoption curves and multi-scenario stochastic demand models rather
than static planning assumptions.

\subsection{Limitations of Traditional Planning Approaches}

Traditional EVCS planning approaches often rely on:
\begin{itemize}
  \item single-objective deterministic optimization (e.g., minimizing
        average distance or total cost alone);
  \item simplified queueing assumptions such as the Erlang-B loss model,
        which ignores physical waiting bays and overestimates blocking;
  \item static demand assumptions that do not account for long-term
        EV adoption growth or stochastic variability;
  \item uniform service radii that ignore urban density heterogeneity
        and road-network accessibility differences.
\end{itemize}
These models are insufficient for city-scale planning where multiple
conflicting objectives must be balanced simultaneously. In particular,
they fail to capture trade-offs between cost, coverage, service quality,
and financial sustainability --- or to enforce service-level constraints
such as maximum blocking probability during peak demand periods.

\subsection{NSGA-II, MOPSO, and Queueing Models}

NSGA-II is an elitist multi-objective genetic algorithm that employs
fast non-dominated sorting with $\mathcal{O}(MN^2)$ complexity and a
crowding distance diversity mechanism. It has become the benchmark
algorithm for multi-objective combinatorial optimisation, with
well-understood convergence guarantees and binary crossover/mutation
operators that map naturally onto station selection problems.

MOPSO extends the Particle Swarm Optimisation paradigm to multiple
objectives through an external archive of non-dominated solutions and
crowding-based leader selection. For binary problems, continuous particle
positions are binarised via a threshold for objective evaluation, while
velocity updates operate in continuous space. MOPSO is known to converge
faster than NSGA-II in early iterations, motivating its inclusion as a
comparative benchmark.

The M/M/c/K finite-capacity queueing model provides analytical
steady-state expressions for queue length $L_q$, waiting time $W_q$,
server utilisation $\rho$, and blocking probability $P_{\text{block}}$.
The finite buffer assumption ($K = c + W$ waiting bays) is more realistic
for urban charging stations than infinite-queue M/M/c models, and
Time-of-Use segmentation of the arrival rate captures the bimodal daily
demand pattern observed in field data.

\section{Motivation}
\label{sec:intro_motivation}

The motivation for this research arises from both practical infrastructure
requirements and methodological gaps in existing EVCS planning literature.

\subsection*{1.~Addressing Infrastructure Deficits in Indian Cities}
Most multi-objective EVCS placement studies have been conducted for
Chinese or European cities using generic urban models. Studies specifically
for Tier-1 Indian cities with real ward-level demographic data, India-specific
policy parameters, and heterogeneous venue types are scarce. Indore, with
its 85~wards, diverse land-use, and active smart-city programme, provides
an empirically grounded and practically relevant case study.

\subsection*{2.~Integrated Four-Objective Formulation}
Existing studies typically address two or three objectives. The simultaneous
optimisation of capital expenditure, spatial coverage, net ROI, and queue
waiting time --- with net ROI defined as annual profit divided by CAPEX,
rather than raw profit, to avoid degeneracy with the coverage objective ---
has not been reported for Indian cities. This formulation creates genuine
trade-off diversity in the Pareto search space.

\subsection*{3.~Realistic Queueing with Blocking Constraints}
Prior work either uses over-simplified Erlang-B loss models or
time-averaged M/M/c queues. An integrated M/M/c/K model with ToU
demand periods and a blocking probability hard constraint embedded
inside the evolutionary evaluation loop represents a methodological
advance that forces the optimizer to provision sufficient charger
capacity at demand bottlenecks.

\subsection*{4.~Stochastic Demand Robustness}
Deterministic demand forecasts are insufficient for infrastructure
planning over a 10-year horizon. Monte~Carlo simulation with 100
demand scenarios ensures that selected station configurations remain
feasible under the full range of plausible future demand trajectories,
not merely under the mean forecast.

\subsection*{5.~Reproducible, Extensible Planning Tool}
The combination of open demographic data, a modular Python
implementation, two well-documented evolutionary algorithms, and
analytical queueing models yields a framework that is reproducible,
extensible to other cities, and capable of accommodating future
enhancements such as renewable integration, vehicle-to-grid (V2G)
capabilities, and real-time dynamic pricing.

\section{Thesis Objectives}
\label{Chapter_1_Objectives}

\noindent The specific objectives of this thesis are:

\begin{enumerate}
  \item To develop a data-driven multi-objective optimization model for
        EVCS placement across Indore's 85~wards, integrating a
        CAGR-based logistic EV adoption curve with Monte~Carlo
        uncertainty quantification to forecast ward-level charging
        demand, and formulating the station selection problem as a
        four-objective binary optimisation incorporating total capital
        expenditure, spatial demand coverage, net return on investment,
        and mean queue waiting time.

  \item To implement and comparatively evaluate NSGA-II and MOPSO on
        the formulated problem, embedding an analytical M/M/c/K
        finite-capacity queueing model with Time-of-Use demand periods
        and a hard blocking probability constraint
        ($P_{\text{block}} \leq 10\%$) directly within the evolutionary
        evaluation loop, and quantifying algorithm quality using
        Hypervolume~(HV), Generational Distance~(GD), and Spread
        performance indicators.

  \item To validate the robustness and scalability of the proposed
        framework through sensitivity analysis with respect to EV
        growth rate, electricity tariff, installation cost, and charger
        service rate, and through scalability benchmarks confirming
        polynomial runtime growth suitable for larger metropolitan
        deployments.
\end{enumerate}

\section{Layout of the Thesis}
\label{sec:intro_layout}

This report is structured into the following chapters:

\begin{itemize}
  \item \textbf{Chapter 2: Literature Review} --- Reviews existing work on
        EVCS placement, multi-objective evolutionary algorithms, M/M/$c$/K
        queueing models, and stochastic demand forecasting methods. It
        identifies five research gaps that motivate the proposed framework
        and provides a structured comparative summary of twelve
        representative studies.

  \item \textbf{Chapter 3: Model Architecture} --- Describes the complete
        architecture of the proposed framework, including the deterministic
        baseline model, four-objective problem formulation, NSGA-II and
        MOPSO optimization layers, M/M/$c$/K queueing evaluation layer with
        Time-of-Use demand segmentation, and formal algorithm performance
        evaluation using HV, GD, and Spread indicators.

  \item \textbf{Chapter 4: Methodology} --- Details the structured
        multi-stage pipeline covering dataset acquisition, spatial zoning,
        ward-level logistic demand forecasting, Monte~Carlo stochastic
        simulation, four-objective optimization using NSGA-II and MOPSO,
        queueing-based operational evaluation, sensitivity analysis, and
        scalability benchmarking.

  \item \textbf{Chapter 5: Results and Analysis} --- Presents Pareto front
        comparisons between NSGA-II and MOPSO, formal algorithm performance
        indicators, optimal solution analysis, queue waiting time
        distributions, HV convergence behaviour, sensitivity studies on
        EV growth rate and tariff parameters, and scalability benchmarks.

  \item \textbf{Chapter 6: Conclusion and Future Work} --- Summarizes the
        main contributions, states deployment recommendations for Indore,
        and outlines future extensions including road-network routing,
        heterogeneous charger-type optimisation, battery swapping station
        integration, and statistical significance testing.
\end{itemize}
% End of Chapter 1